\documentclass[10pt, letterpaper]{article}

\usepackage[margin=0.75in]{geometry}
\usepackage{enumitem}
\usepackage{hyperref}
\usepackage{titlesec}
\usepackage[T1]{fontenc}
\usepackage[utf8]{inputenc}

\hypersetup{
  colorlinks=true,
  urlcolor=black,
  linkcolor=black
}

\setlength{\parindent}{0pt}
\setlength{\parskip}{0pt}

\titleformat{\section}
  {\normalfont\normalsize\bfseries}
  {}{0em}{}
  [\vspace{-4pt}\rule{\linewidth}{0.5pt}\vspace{-6pt}]

\titlespacing{\section}{0pt}{10pt}{6pt}

\setlist[itemize]{
  leftmargin=1.5em,
  itemsep=1pt,
  parsep=0pt,
  topsep=2pt
}

\newcommand{\role}[2]{%
  \textit{#1} \hfill #2 \\[-2pt]%
}

\begin{document}

% ── Header ────────────────────────────────────────────────────────────────────
\begin{center}
  {\Large \textbf{Hithesh Rai Purushothama}} \\[4pt]
  \small
  480-289-0708 \,\textbar\,
  \href{mailto:hraipuru@asu.edu}{hraipuru@asu.edu} \,\textbar\,
  \href{https://linkedin.com/in/hithesh-rai-p}{linkedin.com/in/hithesh-rai-p} \,\textbar\,
  \href{https://github.com/hitheshrai}{github.com/hitheshrai} \,\textbar\,
  \href{https://hitheshrai.github.io/Hithesh}{hitheshrai.github.io/Hithesh}
\end{center}

% ── Professional Summary ──────────────────────────────────────────────────────
\section*{Professional Summary}

Materials and AI engineering graduate student with hands-on research experience spanning
photovoltaic materials, electrochemical impedance spectroscopy, and physics-informed machine
learning across four international research institutions. Builds end-to-end systems across the
full stack: from thin-film device fabrication and EIS characterization to GPR-based degradation
prediction pipelines, edge AI deployment on constrained hardware, and autonomous robotic systems.
Equally comfortable in a materials science lab, an HPC environment, or a robotics integration
workflow. Interested in self-driving laboratory systems, transferable ML frameworks for energy
materials, and deploying AI in low-connectivity or resource-constrained environments.

% ── Education ─────────────────────────────────────────────────────────────────
\section*{Education}

\textbf{Arizona State University} \hfill Tempe, AZ \\
\role{M.S.\ in Artificial Intelligence Engineering (Materials Science and Engineering)}{Jan 2026 -- May 2027}

\vspace{4pt}
\textbf{Arizona State University} \hfill Tempe, AZ \\
\role{B.S.E.\ Electrical and Electronic Engineering}{Aug 2021 -- Dec 2025}

% ── Research Experience ───────────────────────────────────────────────────────
\section*{Research Experience}

% --- Rolston Lab ---
\textbf{Renewable Energy Materials and Devices Lab, Arizona State University} \hfill Tempe, AZ \\
\role{Research Assistant}{Sep 2022 -- Present}
\begin{itemize}
  \item Developing an ML pipeline using Gaussian Process Regression (GPR) on electrochemical
    impedance spectroscopy (EIS) features to predict battery capacity fade and interfacial
    degradation, with application to state-of-health estimation across energy storage systems.
  \item Reproduced and extended Zhang et al.\ (2020, \textit{Nature Communications}) battery
    degradation model, achieving $R^2 = 0.81$ on unseen test cells; identified a 17.80\,Hz EIS
    feature as a degradation-robust predictor across temperature regimes using ARD kernel
    weighting.
  \item Applying multitask and transfer learning to model interfacial degradation across batteries,
    perovskite solar cells, and fuel-cell-relevant material systems from shared EIS descriptors.
  \item Automated EIS data preprocessing, Distribution of Relaxation Times (DRT) peak extraction,
    and GPR model training pipeline in Python (NumPy, GPyTorch, scikit-learn) on ASU's Sol HPC
    cluster via VSCode remote integration.
  \item Fabricated and optimized perovskite thin films using spin coating and blade coating on
    rigid and flexible substrates, achieving uniform films with reduced defect density.
  \item Conducted optical and electrical characterization with solar simulator, XRD, and
    profilometry to analyze degradation and identify key failure mechanisms guiding material
    improvements.
  \item Presented work at IEEE PVSC 2024 (Seattle) and IPERO 2025 (Kyoto); submitted abstract to
    AI4X -- Accelerate Conference 2026 (Singapore) with advisor N.\ Rolston.
\end{itemize}

\vspace{6pt}

% --- EPFL ---
\textbf{Photovoltaics Lab, \'Ecole Polytechnique F\'ed\'erale de Lausanne (EPFL)} \hfill Neuch\^atel, Switzerland \\
\role{Undergraduate Research Assistant --- ThinkSwiss Scholar}{May 2025 -- Aug 2025}
\begin{itemize}
  \item Fabricated single-junction perovskite thin-film devices achieving efficiencies up to 19\%
    using controlled spin-coating workflows with systematic process parameter variation.
  \item Deposited SnO$_2$ electron transport layers via ALD and analyzed thickness and conformity;
    used thermal evaporation for transport layers and metal contacts, improving device yield.
  \item Assisted with encapsulation and stability testing, identifying key degradation pathways
    and informing process-improvement actions.
\end{itemize}

\vspace{6pt}

% --- HZB ---
\textbf{Helmholtz Zentrum Berlin --- Quantum Materials Group} \hfill Berlin, Germany \\
\role{International Summer Student}{Jul 2024 -- Aug 2024}
\begin{itemize}
  \item Investigated structural instabilities in perovskite-based ferroelectric and piezoelectric
    systems using Pair Distribution Function (PDF) analysis of X-ray and neutron diffraction data.
  \item Identified composition-driven morphotropic phase transitions and correlated atomic-scale
    structural variations with compositional changes using Python analysis scripts.
\end{itemize}

\vspace{6pt}

% --- Purdue ---
\textbf{Letian Dou Group, Purdue University} \hfill West Lafayette, IN \\
\role{Summer Undergraduate Research Fellow}{May 2023 -- Aug 2023}
\begin{itemize}
  \item Supported halide perovskite stability studies through device fabrication and testing;
    reviewed and organized five years of literature on additive engineering and degradation
    mechanisms.
  \item Compiled a comparative Python database of device structures, processing conditions, and
    efficiencies; used the Perovskite Database ML platform to identify material-performance trends
    guiding additive selection.
\end{itemize}

% ── Engineering Experience ────────────────────────────────────────────────────
\section*{Engineering Experience}

\textbf{Capstone Project --- Power System Modeling for Grid-Scale Energy Storage} \hfill Tempe, AZ \\
\role{Arizona State University}{Jan 2025 -- Present}
\begin{itemize}
  \item Modeled Battery Energy Storage Systems (BESS) for renewable energy integration using
    PyPSA, Pandas, and NumPy, producing performance summaries that informed storage sizing and
    operational trade-off decisions.
  \item Deployed an LLM-powered digital twin dashboard to monitor and interpret real-time PV +
    BESS telemetry, enabling early fault detection and interpretable system diagnostics.
  \item Designed realistic grid-scale load profiles from utility and consumer datasets to evaluate
    system-level PV and storage performance.
\end{itemize}

% ── Leadership \& Management ──────────────────────────────────────────────────
\section*{Leadership \& Management}

% --- Management Intern ---
\textbf{Next Lab, Arizona State University} \hfill Tempe, AZ \\
\role{Management Intern}{Mar 2026 -- Present}
\begin{itemize}
  \item Lead evaluation and selection of emerging AI technologies for lab adoption, assessing
    technical maturity, integration feasibility, and deployment cost-benefit across candidate
    systems --- serving as the lab's primary technology intelligence function.
  \item Drive end-to-end AI project execution from scoping through implementation, coordinating
    across research and operations staff to deliver production-ready systems on schedule.
  \item Support core lab leadership on strategic initiatives, translating complex technical
    findings into actionable recommendations for non-technical stakeholders.
\end{itemize}

\vspace{6pt}

% --- Senior Studio Associate ---
\textbf{Next Lab, Arizona State University} \hfill Tempe, AZ \\
\role{Senior Studio Associate --- AI \& Edge Systems}{Sep 2023 -- Present}
\begin{itemize}
  \item Built end-to-end benchmarking pipelines in Python to evaluate AI inference systems across
    latency, memory, accuracy, hallucination rate, and power consumption on NVIDIA Jetson devices,
    informing deployment decisions for production edge environments.
  \item Analyzed trade-offs between INT8/Q4 quantization strategies and model performance,
    identifying failure modes and guiding optimization for robust real-world deployment at scale.
  \item Designed and validated agent-based task execution and knowledge-retrieval pipelines using
    LangChain, reducing manual workflow overhead and improving automation reliability.
  \item Deployed a digital twin dashboard powered by an LLM to monitor and interpret PV + BESS
    system telemetry in real time, enabling early fault detection across operational scenarios.
\end{itemize}

\vspace{6pt}

% --- EPICS ---
\textbf{Engineering Projects in Community Service (EPICS), Arizona State University} \hfill Tempe, AZ \\
\role{Team Lead}{Jan 2022 -- Jan 2024}
\begin{itemize}
  \item Led a team of 10+ students to develop a Python and robotics curriculum for middle
    schoolers across the US and South Africa, delivering hands-on embedded AI modules using
    NVIDIA Jetson hardware with Jetson Infinity.
  \item Applied NumPy and Matplotlib for simulation components; enhanced learning outcomes through
    robotic arm hands-on activities.
\end{itemize}

% ── Selected Projects ─────────────────────────────────────────────────────────
\section*{Selected Projects}

\textbf{NEXUS --- Automated Research Intelligence System} \hfill Next Lab, ASU
\begin{itemize}
  \item Built an automated morning briefing pipeline delivering daily AI/energy research summaries
    to WhatsApp using local LLMs on an NVIDIA DGX Spark, with a rolling context file system and
    automated references pipeline for citation-grounded output.
  \item Deployed and debugged NVIDIA NIM containers on DGX Spark, resolving permission, Docker
    socket, and NGC authentication issues; implemented Tailscale for secure remote access.
\end{itemize}

\vspace{6pt}

\textbf{Nex Bot --- Autonomous Voice Receptionist} \hfill Next Lab, ASU
\begin{itemize}
  \item Deployed Reachy Mini humanoid robot as an AI-powered voice receptionist with full
    inference offloaded to a DGX Spark over Tailscale, enabling low-latency on-device interaction
    without cloud dependency.
  \item Integrated faster-whisper (STT), Silero VAD, Kokoro TTS, and a local LLM into a
    real-time conversational pipeline using the \texttt{reachy-mini} PyPI stack.
\end{itemize}

\vspace{6pt}

\textbf{Autonomous Social Robot System} \hfill Next Lab, ASU
\begin{itemize}
  \item Built a fully offline public-facing autonomous robot using Reachy Mini + DGX Spark,
    integrating YOLOv8n for person detection, Silero VAD for voice activity detection, and a
    complete STT $\to$ LLM $\to$ TTS pipeline running without internet connectivity.
  \item Resolved Zenoh daemon connectivity and Docker stack configuration issues to achieve
    stable real-time inference across all modalities.
\end{itemize}

\vspace{6pt}

\textbf{EIS-GPR Battery Degradation Pipeline} \hfill Rolston Lab, ASU
\begin{itemize}
  \item Reproduced Zhang et al.\ (2020) in both scikit-learn and GPyTorch; extended to
    multi-temperature training with ARD kernels, isolating SEI-layer feature at 17.80\,Hz as a
    cross-temperature degradation predictor ($R^2 = 0.81$ on held-out test cells).
  \item Documented full methodology in a CLAUDE.md session file; planning Fig.\ 1--2 reproduction
    using 25$^\circ$C test cells and extension to DRT-derived features.
\end{itemize}

% ── Conferences ───────────────────────────────────────────────────────────────
\section*{Conferences and Presentations}

\begin{itemize}
  \item AI4X -- Accelerate Conference 2026, Singapore (accepted) ---
    ``Transferable Impedance-Grounded Learning for Interfacial Degradation Across Energy Systems''
    (with N.\ Rolston).
  \item IEEE Photovoltaic Specialists Conference, June 2024, Seattle, WA ---
    ``Open-Air Processing of Cesium-Based Inorganic Perovskites for Radioisotopes.''
  \item International Conference on Perovskite, Organic Photovoltaics and Optoelectronics,
    Jan 2025, Kyoto, Japan ---
    ``Blade-Coated Cesium Lead Halide Perovskite Thin Films for Alphavoltaic Applications.''
  \item Stanford Research Conference, April 2024, San Jose, CA ---
    ``Enhancing the Performance of Formamidinium Lead Halide Perovskite Solar Cells.''
  \item Arizona Student Energy Conference (AzSEC), April 2023, Tempe, AZ ---
    ``Unlocking the Potential of Halide Perovskites: The Art of Thin Film Fabrication.''
\end{itemize}

% ── Awards ────────────────────────────────────────────────────────────────────
\section*{Awards}

\begin{itemize}
  \item \textbf{ThinkSwiss Research Scholarship}, Swiss National Science Foundation \hfill Feb 2025 \\
    \small Awarded for summer research at EPFL on perovskite thin-film device fabrication and
    stability studies. \hfill Washington, DC
  \item \textbf{Fulton Undergraduate Research Initiative Grant (FURI)} \hfill Fall/Spring 2023--2025 \\
    \small Arizona State University, Tempe, AZ
  \item \textbf{Purdue Summer Undergraduate Research Fellowship} \hfill Summer 2023 \\
    \small Purdue University, West Lafayette, IN
\end{itemize}

% ── Technical Skills ──────────────────────────────────────────────────────────
\section*{Technical Skills}

\begin{itemize}
  \item \textbf{ML \& Data:} Gaussian Process Regression, Transfer/Multitask Learning,
    Time-Series Forecasting, Bayesian Modeling, scikit-learn, GPyTorch, TensorFlow, NumPy, Pandas
  \item \textbf{Energy Systems:} EIS, DRT Analysis, PAIOS, PyPSA, PVlib, GridLAB-D,
    MATLAB/Simulink, BESS Modeling, Load Flow Analysis, SCADA/EMS Basics
  \item \textbf{Fabrication:} Spin Coating, Blade Coating, ALD, Thermal Evaporation,
    Encapsulation
  \item \textbf{Characterization:} Solar Simulator, XRD, SEM, Profilometry, Optical Inspection,
    Pair Distribution Function (PDF) Analysis
  \item \textbf{Edge AI \& Inference:} ONNX Runtime, TensorRT, MLC, INT8/Q4 Quantization,
    NVIDIA Jetson, NVIDIA DGX Spark, YOLOv8n
  \item \textbf{Robotics \& Autonomy:} Reachy Mini (\texttt{reachy-mini}), faster-whisper,
    Silero VAD, Kokoro TTS, Zenoh, LangChain
  \item \textbf{Infrastructure:} Docker, GitHub, Tailscale, Embedded Linux (Jetson/Pi/DGX),
    HPC (Sol/SLURM), VSCode Remote
  \item \textbf{Programming:} Python, C++, MATLAB, Verilog
  \item \textbf{EDA Tools:} ModelSim, Quartus
\end{itemize}

\end{document}